\BookPart{跋}
丙午之世，屏间奇谈，无远弗届。有人据一字而定万国之疆，有人依一画而辨千年之史；有人以史为镜，有人以镜为史。其说虽未必有据，其传则确有迹。故备录于此，非以证往古，乃以存今人。

倘后世览之，宜辨其辞而究其据，勿以声音相似便论渊源，勿以疑心深重便为证据，尤勿以人群之异而断其为非人。史学贵考实，求知贵存疑；能信古而不泥古，能爱国而不造敌，庶几不负寻根之本意。

是故：

\vspace{8mm}
\begin{verseblock}
丙午奇谈遍九州，方屏咫尺定春秋。 \\
千年帝系随音转，万里山河逐字收。 \\
有据偏疑书尽伪，无凭却道世同谋。 \\
他年若问寻根事，一卷存疑待考求。
\end{verseblock}
\begin{flushright}\vspace{9mm}\textcolor{wine}{丙午岁，编者谨识。}\end{flushright}
